\section{执行、做市与库存控制}
\label{ch:lower-microstructure-control}

\begin{itemize}
  % Generated from curriculum/manifest.json; do not edit by hand.
\item 直接先修：下册《事件与订单簿：从条件强度到队列成交》、上册第 13 章《凸优化、KKT 与对偶》、上册第 14 章《浮点数、数值线代与病态问题》、上册第 15 章《Monte Carlo、bootstrap 与方差缩减》。

\end{itemize}

立即成交会付出冲击，等待又承担价格与未完成风险。将剩余数量作为状态，可以把这两种代价写在同一个决策方程里。


\MFQLead{先定义实施差额的基准}

对买入母单，决策时到达价为 $S_0$，第 $k$ 期实际成交 $x_k$、成交价 $P_k$，实现差额为
\[
 \mathrm{IS}=\sum_k x_k(P_k-S_0)+R_K(S_K-S_0)+C,
\]
其中 $R_K$ 是终止时未完成数量，$C$ 是显式费用。若只统计已成交部分，应明确这是 partial-fill IS，不能把未完成任务静默删除。卖单符号相反，所有价格和费用必须统一为现金或基点口径。

教学冲击模型写成
\[
 P_k=S_0+\gamma X_{k-1}+\eta x_k,
\]
其中 $X_{k-1}=\sum_{j<k}x_j$，$\eta x_k$ 是本期临时冲击，$\gamma X_{k-1}$ 是累计永久冲击。永久冲击改变后续基准，临时冲击只改变当前成交价；把二者合并成一个系数会扭曲日程的跨期权衡\cite{almgren2001optimal}。

事先确定计划 $(3,3,3)$ 在完成率 $(1,0.5,0)$ 下实际成交 $(3,2,0)$。以 $S_0=100,\eta=0.2,\gamma=0.1$ 计算，partial-fill IS 为 $3.2$，剩余 4 手必须保留在状态中。

\MFQLead{随机完成率与停止规则}

设计划量 $u_k$，随机完成率 $F_k\in[0,1]$，实际成交为
\[
 x_k=\min\{R_k,\operatorname{round}(u_kF_k)\},
 \qquad R_{k+1}=R_k-x_k.
\]
完成率不是预测误差的装饰项：它决定剩余库存、后续冲击和最终违约状态。停止规则可以是固定期限、累计损失阈值、市场状态或风险限制；一旦触发，应输出 stopped、剩余数量和采用的处置规则。

动态规划的一般值函数记为 $V_t(x)$，表示时点 $t$、剩余库存 $x$ 下的最小未来目标；离散状态取 $(k,R_k)$。一个最小成本递推是
\[
 V_k(R)=\min_{0\le x\le\min(R,x_{\max})}
 \left\{\eta x^2+\gamma(I-R)x+\phi(R-x)^2+V_{k+1}(R-x)\right\},
\]
终端只允许 $R=0$。事先确定六手、三期问题的最优日程是 $(3,2,1)$，目标 $19.2$；完整枚举所有可行路径得到相同答案。DP 证明的是离散目标的最优性，不证明冲击参数、二次风险或硬完成约束适合现实委托。

\MFQLead{库存如何影响下一次报价}

令中价 $m_t$、半点差 $h$、库存 $q_t$、库存偏移系数 $\kappa$。最小双边报价为
\[
 r_t=m_t-\kappa q_t,
 \qquad b_t=r_t-h,\quad a_t=r_t+h.
\]
若 bid 成交，$q_{t+1}=q_t+1$、现金减少 $b_t$；若 ask 成交，库存减一、现金增加 $a_t$。下一轮必须用 $q_{t+1}$ 计算报价。算例从零库存开始，首轮 bid 成交后库存为 1，下一轮报价从 $99.5/100.5$ 下移为 $99.4/100.4$。双边同时下移，使买入变得不积极、卖出更容易，才形成真正反馈。

% Generated by tools/build_figures.py; edit figures/figure-specs.json instead.
\MFQFigure{tex/figures/generated/lower-ms-feedback.pdf}{每次成交改变现金与库存，新的库存影响下一次报价；库存限制也会改变后续可选动作。}{lower-ms-feedback}


成交改变库存，新的库存又影响下一次报价，这构成动态反馈。如果预先生成与库存无关的全天报价，再计算库存路径，研究的是另一种不含库存反馈的策略；如果报价使用了尚未发生的库存或事件，则超出了当时的信息。进一步的模型可以考虑成交强度不确定性、价格跳跃、库存边界和撤单延迟。

\MFQTransition{Almgren--Chriss 的风险--冲击折中}

离散执行把库存 $x_k$ 从 $X$ 降到 0，交易量 $n_k=x_{k-1}-x_k$。典型目标为期望成本加方差惩罚：
\[
\min_{n_1,\ldots,n_K}
\mathbb E[C]+\lambda\operatorname{Var}(C),
\qquad \sum_kn_k=X.
\]
临时冲击影响本次成交，永久冲击改变后续基准。$\lambda=0$ 更愿意等待低冲击，$\lambda$ 大则加快执行以降低价格风险。参数应从订单级数据估计并做容量压力，不能从公开日频量价唯一识别。

% Generated by tools/build_figures.py; edit figures/figure-specs.json instead.
\MFQFigure{tex/figures/generated/lower-ms-control.pdf}{执行与做市控制根据库存和剩余时间选择动作，再由随机成交改变下一状态；终端惩罚约束未完成风险。}{lower-ms-control}


参与率、最小成交单位、最大切片和收盘前强制完成使问题成为约束动态规划。未完成状态必须在终端施加机会成本或硬约束，否则最优策略可能选择永不交易。

\MFQTransition{做市的库存控制}

简化 Avellaneda--Stoikov 中，库存 $q$ 改变保留价格
\[
r_t=S_t-q\gamma\sigma^2(T-t).
\]
多头库存使报价中心下移，鼓励卖出；风险厌恶、波动和剩余时间放大调整。最优价差还取决于订单到达对距离的敏感度。模型假设连续 Brownian 中价和指数到达，真实跳跃、队列与延迟都可能破坏闭式。

报价状态机还要处理库存上限、单边停报、撤单确认、自成交防护和 stale quote。控制动作必须在收到事件后生成，并等待交易所确认；本地认为已撤单不代表市场上已撤。

\MFQTransition{部分可观测与鲁棒控制}

真实状态包含不可见的短期 alpha、对手意图和未来流动性，执行问题更接近 POMDP。滤波器提供 belief state，控制器基于 belief 行动。若状态估计和控制用同一未来数据平滑，会产生前视。

模型参数不确定时可做情景 DP、分布鲁棒或保守策略。最简单的鲁棒基线是限制参与率、库存和单步动作，并在成交率低于阈值时切换预定义应急日程。复杂 RL 只有在能恢复小型 DP 解析参照 后才值得使用。

\begin{diagnostic}
控制器的目标函数必须包含未完成、库存和成本。若奖励只含已实现价差，代理可以通过不成交或积累巨大库存获得虚假高分。
\end{diagnostic}

控制规则将在微观结构路线综合项目中与事件仿真、执行成本和库存压力一起复核。
\subsection{两期执行的风险权衡}

需要买入总量 $Q$，第一期买 $x$，第二期买 $Q-x$，暂设均可成交。临时冲击成本为 $\eta[x^2+(Q-x)^2]$，持有未完成数量的风险惩罚为 $\phi(Q-x)^2$。对 $x$ 求导得
\[
 2\eta x-2(\eta+\phi)(Q-x)=0,
 \qquad x^*=\frac{\eta+\phi}{2\eta+\phi}Q.
\]
$\phi=0$ 时均分；$\phi$ 增加时提前买入更多。$Q=6,\eta=1,\phi=2$ 时 $x^*=4.5$，第二期为 1.5；若只能买整数，比较 $x=4,5$，两者目标都为 28，形成并列最优。

报价反馈也可直接算现金。中价 100、半价差 0.5，从零库存买入一股于 99.5，现金为 $-99.5$、库存为 1，盯市财富为 0.5。若中价随即跌到 99，财富变成 $-0.5$；成交并不保证赚钱，逆向选择可以超过点差收益。notebook 枚举整数日程并显示每次成交后的库存与现金。
